% page 137 du fac-simile (image p-136 du scan)
% bloc 167.6 x 254.9 mm ; marge gauche 34.4 mm
\pgc{12.700mm}{0.000mm}{
\l{}
\l{}
\l{\cel{54}127}
\l{}
\l{}
\l{eglefino.- (zool.) Fisho de la genero \textquotedbl{}gadi\textquotedbl{}. - L.\cel{1}\sou{gadus}}
\l{\cel{3}aeglefinus. - FL.}
\l{}
\l{\sou{ego}. - To ye quo konsistas la individueso, la personeso. - L.}
\l{}
\l{\sou{egorjar}. - (trans.) Ocidar per kultelo-frapo en la kolo (fauco).}
\l{\cel{3}- F.}
\l{}
\l{\sou{egotismo}.- Manio parolar pri su. - EFIRS.}
\l{}
\l{\sou{egotisto}.- Maniiko pri parolar pri su. - DEFIRS.}
\l{}
\l{\sou{\textquotedbl{}eis\textquotedbl{}}. -\cel{2}(\sou{muziko})\cel{2}La noto triesma di la \textquotedbl{}do\textquotedbl{}-gamo diezoza.}
\l{}
\l{\sou{ejektar}. - (\sou{trans.}) I. (\sou{mekaniko}). Retrojetar adextere aero per}
\l{\cel{3}aparato qua evakuas ta fluido per altra fluid-amaso qua movas}
\l{\cel{3}tre rapide. - II. (\sou{Pafarmo}). Ekjetar (la mufo di kartocho).}
\l{\cel{3}- DEFIS.}
\l{}
\l{ek. - I. de interne ad extere... - II. Indikas la materio di kozo}
\l{\cel{3}(de qua onu imaginas ke lu esas extraktita). - III. (\sou{metaf.})}
\l{\cel{3}Indikas objekto, o mem ento, qua apartenas ad ensemblo,}
\l{\cel{3}kolektitaro (e quin onu supozas extraktita de olu).}
\l{}
\l{\sou{eko}. (\sou{Akust}.) Reflekto de la son-ondi retrosendita a la oreli}
\l{\cel{3}da surfaco qua produktas retroshoko. - DEFIRS.}
\l{}
\l{\sou{\textquotedbl{}ekarte\textquotedbl{}}.- Kartoludo quan partoprenas du personi, e segun la}
\l{\cel{3}regulo di qua la ludanti darfas forjetar omna o kelka karti,}
\l{\cel{3}e prenar altrai kom remplasanta. - F.}
\l{}
\l{\sou{ekidno}.- (\sou{zool.})\cel{2}Mamifero Australiana, de la familio \textquotedbl{}sen-denti\textquotedbl{},}
\l{\cel{3}di qua la muzelo esas longa e tenua, la korpo kovrata da pik-}
\l{\cel{3}anti quale la herisono, qua ovifas e nutras su ek insekti. -}
\l{\cel{3}L.\cel{1}\sou{echidna (hystrix)}. - DEFIS.}
\l{}
\l{\sou{ekimoso}.- (\sou{patol.})\cel{3}Makulo livida sur la pelo, produktita da}
\l{\cel{3}la ekfluo di la sango en la tisuo celulara. - DEFIS.}
\l{}
\l{\sou{ekino}.- (\sou{zool.}). Ekinodermo globatra, herisata ye pikanti (no-}
\l{\cel{3}mata anke \textquotedbl{}mar-herisono\textquotedbl{}.) - L.\cel{1}\sou{echinus}.}
\l{}
\l{\sou{ekinodermo}.- (\sou{zool.}) Klaso de animali radioza, di qui la korpo}
\l{\cel{3}esas garnisita ye multega pinti artikoza e movebla.- DEFIS.}
\l{}
\l{\sou{eklampsio}.- (\sou{patol.}) Konvulsi qui eventas maxim-multa-kaze pro}
\l{\cel{3}albuminurio. Konvulsi che infanteti. - DEFIS.}
\l{}
\l{\sou{eklektika}.- Qua asemblas en un doktrino-korpo la exaktaji quaze}
\l{\cel{3}dissemita en sistemi diversa di filozofio e di cienco. - DEFIRS}
\l{}
\l{\sou{eklezio}.- Asemblitaro de ti qui adoras la deo di la kristani :}
\l{\cel{3}la eklezio kristana, asemblitaro de ti qui profesas la lego di}
\l{\cel{3}Kristo. - EFIRS.}
}
